\documentclass[10pt]{report}

% ============================================================
% PAGE SETUP
% ============================================================
\usepackage[a4paper, margin=1in]{geometry}

% ============================================================
% CORE PACKAGES
% ============================================================
\usepackage[T1]{fontenc}
\usepackage[utf8]{inputenc}
\usepackage[australian]{babel}
\usepackage{textcomp}

\usepackage{graphicx}
\usepackage{float}
\usepackage{booktabs}
\usepackage{array}
\usepackage{tabularx}
\usepackage{makecell}
\usepackage{longtable}

\usepackage{amsmath}
\usepackage{mathtools}

\usepackage{tikz}
\usepackage{xcolor}
\usepackage[table]{xcolor}

\usepackage{caption}
\usepackage{fancyhdr}
\usepackage{titlesec}
\usepackage{listings}
\usepackage{needspace}

% ============================================================
% HYPERLINKS
% ============================================================
\usepackage[
    colorlinks=true,
    linkcolor=blue,
    citecolor=blue,
    urlcolor=blue
]{hyperref}

% ============================================================
% SECTION FORMATTING
% ============================================================
\renewcommand{\thesection}{\Roman{section}}
\renewcommand{\thesubsection}{\alph{subsection}.}

\titleformat{\section}
  {\centering\normalfont\Large\bfseries}
  {\thesection.}
  {0.5em}
  {}

\titleformat{\subsection}
  {\normalfont\bfseries\normalsize}
  {\thesubsection}
  {0.5em}
  {}

\titlespacing*{\section}
{0pt}{1.0ex plus 0.2ex minus 0.2ex}{0.8ex}

\titlespacing*{\subsection}
{0pt}{0.8ex plus 0.2ex minus 0.2ex}{0.5ex}

% ============================================================
% CAPTION FORMATTING
% ============================================================
\captionsetup[table]{
    position=above,
    justification=centering,
    singlelinecheck=false,
    font=small,
    labelfont=bf,
    textfont=it,
    labelsep=period,
    skip=4pt
}

\captionsetup[figure]{
    position=below,
    justification=centering,
    singlelinecheck=false,
    font=small,
    labelfont=bf,
    textfont=it,
    labelsep=period,
    skip=4pt
}

\captionsetup[lstlisting]{
    justification=raggedright,
    singlelinecheck=false,
    font=small,
    labelfont=bf,
    skip=2pt
}

% ============================================================
% HEADER / FOOTER
% ============================================================
\pagestyle{fancy}
\fancyhf{}
\fancyhead[L]{C++ Programming Notes}
\fancyhead[R]{Personal Learning}
\fancyfoot[C]{C++ Language Learning Notes}
\fancyfoot[R]{\thepage}

% ============================================================
% COLOURS
% ============================================================
\definecolor{CoverBlue}{RGB}{0,32,96}
\definecolor{codebg}{RGB}{245,245,244}
\definecolor{codeframe}{RGB}{180,180,180}
\definecolor{codeblue}{RGB}{0,70,180}
\definecolor{codegreen}{RGB}{0,120,60}
\definecolor{codepurple}{RGB}{120,40,140}
\definecolor{codenumbers}{RGB}{120,120,120}
\definecolor{codered}{RGB}{160,30,30}

% ============================================================
% C++ CODE STYLE
% ============================================================
\lstdefinestyle{CppStyle}{
    language=C++,
    basicstyle=\ttfamily\small,
    backgroundcolor=\color{codebg},
    frame=single,
    rulecolor=\color{codeframe},
    framerule=0.25pt,
    framesep=4pt,
    breaklines=true,
    breakatwhitespace=false,
    showstringspaces=false,
    tabsize=4,
    numbers=left,
    numberstyle=\tiny\color{codenumbers},
    numbersep=6pt,
    xleftmargin=0pt,
    xrightmargin=0pt,
    framexleftmargin=18pt,
    framexrightmargin=4pt,
    framextopmargin=4pt,
    framexbottommargin=4pt,
    keywordstyle=\color{codeblue}\bfseries,
    commentstyle=\color{codegreen}\itshape,
    stringstyle=\color{codered},
    identifierstyle=\color{black},
    captionpos=t,
    aboveskip=6pt,
    belowskip=6pt,
    linewidth=\linewidth,
    columns=fullflexible,
    keepspaces=true
}

\lstset{style=CppStyle}

% ============================================================
% CUSTOM NOTE BOXES
% ============================================================
\newcommand{\conceptbox}[1]{
    \vspace{0.5em}
    \noindent
    \fcolorbox{CoverBlue}{blue!5}{
        \begin{minipage}{0.96\linewidth}
        #1
        \end{minipage}
    }
    \vspace{0.5em}
}

\newcommand{\important}[1]{
    \vspace{0.5em}
    \noindent
    \fcolorbox{red!60!black}{red!5}{
        \begin{minipage}{0.96\linewidth}
        \textbf{Important:} #1
        \end{minipage}
    }
    \vspace{0.5em}
}

\newcommand{\examplebox}[1]{
    \vspace{0.5em}
    \noindent
    \fcolorbox{CoverBlue}{blue!4}{
        \begin{minipage}{0.96\linewidth}
        \textbf{Example:} #1
        \end{minipage}
    }
    \vspace{0.5em}
}

% ============================================================
% DOCUMENT START
% ============================================================
\begin{document}

% ============================================================
% TITLE PAGE
% ============================================================
\begin{titlepage}
\thispagestyle{empty}

\begin{tikzpicture}[remember picture,overlay]
    \draw[CoverBlue, line width=1.2pt]
        ([xshift=1.3cm,yshift=-1.3cm]current page.north west)
        rectangle
        ([xshift=-1.3cm,yshift=1.3cm]current page.south east);
\end{tikzpicture}

\centering

\vspace*{1.2cm}

% ------------------------------------------------------------
% TITLE BLOCK
% ------------------------------------------------------------
{\fontsize{30}{36}\selectfont\bfseries
C++ Programming Language
\par}

\vspace{0.35cm}

{\fontsize{19}{23}\selectfont\bfseries
Personal Learning Notes
\par}

\vfill

% ------------------------------------------------------------
% DESCRIPTION BLOCK
% ------------------------------------------------------------
\begin{minipage}{0.78\textwidth}
\centering
{\fontsize{13}{17}\selectfont\itshape
A structured notebook for learning C++ syntax, program structure,
compilation, functions, data types, memory concepts, object-oriented
programming, and applied software engineering.
\par}
\end{minipage}

\vfill

% ------------------------------------------------------------
% AUTHOR BLOCK
% ------------------------------------------------------------
{\fontsize{14}{18}\selectfont
Prepared by
\par}

\vspace{0.35cm}

{\fontsize{23}{28}\selectfont\bfseries
Chabod Masere
\par}

\vspace{0.6cm}

\vfill

% ------------------------------------------------------------
% DOCUMENT TYPE BLOCK
% ------------------------------------------------------------
{\fontsize{13}{17}\selectfont\bfseries
Personal Study Notes
\par}

\vspace{0.35cm}

{\fontsize{13}{17}\selectfont
\today
\par}

\vspace*{1.2cm}

\end{titlepage}

\clearpage

% ============================================================
% TABLE OF CONTENTS
% ============================================================
\tableofcontents
\clearpage

% ============================================================
% SECTION I
% ============================================================
\section{Introduction to C++}

\subsection{What is C++?}

C++ is a general-purpose programming language widely used for systems programming, embedded software, robotics, simulation, game engines, high-performance computing, and engineering applications.

It provides low-level control over memory and hardware while also supporting high-level programming techniques such as object-oriented programming, generic programming, and abstraction.

\conceptbox{
C++ is powerful because it allows direct control over computation, memory, and performance. This makes it useful for engineering systems where speed, reliability, and hardware-level control matter.
}

\subsection{Why Learn C++?}

C++ is valuable because it develops a strong understanding of:

\begin{itemize}
    \item Program structure
    \item Compilation and linking
    \item Memory management
    \item Data types
    \item Functions
    \item Object-oriented programming
    \item Performance-aware software design
    \item Embedded and robotics software foundations
\end{itemize}

\important{
C++ gives the programmer significant control, but that control also creates responsibility. Poor memory management, invalid pointers, and unsafe assumptions can create serious bugs.
}

% ============================================================
% SECTION II
% ============================================================
\section{Basic Program Structure}

\subsection{Hello World}

\begin{lstlisting}[caption={Hello World in C++}]
#include <iostream>

int main() {
    std::cout << "Hello, World!" << std::endl;
    return 0;
}
\end{lstlisting}

\subsection{Program Explanation}

The basic C++ program contains several important parts:

\begin{itemize}
    \item \texttt{\#include <iostream>} includes the standard input-output stream library.
    \item \texttt{int main()} defines the entry point of the program.
    \item \texttt{std::cout} sends output to the terminal.
    \item \texttt{<<} is the stream insertion operator.
    \item \texttt{std::endl} inserts a newline and flushes the output stream.
    \item \texttt{return 0;} indicates successful program execution.
\end{itemize}

\examplebox{
Every standard C++ program begins execution from the \texttt{main()} function.
}

% ============================================================
% SECTION III
% ============================================================
\section{Compilation and Execution}

\subsection{The Development Cycle}

The C++ development cycle usually follows this process:

\begin{enumerate}
    \item Write the source code.
    \item Compile the source file.
    \item Link object files and libraries.
    \item Produce an executable.
    \item Run the program.
    \item Debug and improve the code.
\end{enumerate}

\conceptbox{
Compilation translates human-readable C++ source code into machine-level instructions that the computer can execute.
}

\subsection{Compiling on Windows}

Assume the source file is named \texttt{hello\_world.cpp}.

\begin{lstlisting}[caption={Compile C++ on Windows}]
g++ hello_world.cpp -o hello_world.exe
\end{lstlisting}

Run the executable:

\begin{lstlisting}[caption={Run Executable on Windows}]
.\hello_world.exe
\end{lstlisting}

\subsection{Compiling on macOS or Linux}

\begin{lstlisting}[caption={Compile C++ on macOS or Linux}]
g++ hello_world.cpp -o hello_world
\end{lstlisting}

Run the executable:

\begin{lstlisting}[caption={Run Executable on macOS or Linux}]
./hello_world
\end{lstlisting}

\important{
Windows commonly uses the \texttt{.exe} executable extension. macOS and Linux executables typically do not require this extension.
}

% ============================================================
% SECTION IV
% ============================================================
\section{Preprocessor Directives and Headers}

\subsection{The Include Directive}

The line:

\begin{lstlisting}[caption={Include Directive}]
#include <iostream>
\end{lstlisting}

instructs the preprocessor to include the standard input-output stream library before compilation.

This allows the program to use objects such as:

\begin{itemize}
    \item \texttt{std::cout}
    \item \texttt{std::cin}
    \item \texttt{std::endl}
\end{itemize}

\subsection{Header Files}

Header files provide declarations for functions, classes, constants, and objects used by the program.

Common C++ headers include:

\begin{table}[H]
\centering
\caption{Common C++ Header Files}
\begin{tabular}{ll}
\toprule
Header & Purpose \\
\midrule
\texttt{<iostream>} & Console input and output \\
\texttt{<string>} & String objects \\
\texttt{<vector>} & Dynamic arrays \\
\texttt{<cmath>} & Mathematical functions \\
\texttt{<fstream>} & File input and output \\
\texttt{<algorithm>} & Common algorithms \\
\bottomrule
\end{tabular}
\end{table}

% ============================================================
% SECTION V
% ============================================================
\section{Namespaces and Output}

\subsection{The Standard Namespace}

The C++ standard library places many features inside the \texttt{std} namespace.

\begin{lstlisting}[caption={Using std Namespace Explicitly}]
#include <iostream>

int main() {
    std::cout << "Learning C++" << std::endl;
    return 0;
}
\end{lstlisting}

\subsection{Scope Resolution Operator}

The symbol \texttt{::} is the scope resolution operator. It tells the compiler where an identifier belongs.

For example:

\begin{lstlisting}[caption={Scope Resolution Operator}]
std::cout
std::endl
\end{lstlisting}

means that \texttt{cout} and \texttt{endl} belong to the \texttt{std} namespace.

\subsection{Using Namespace std}

\begin{lstlisting}[caption={Using Namespace std}]
#include <iostream>

using namespace std;

int main() {
    cout << "Learning C++" << endl;
    return 0;
}
\end{lstlisting}

\important{
Using \texttt{std::} explicitly is usually better in larger projects because it avoids namespace collisions.
}

% ============================================================
% SECTION VI
% ============================================================
\section{Variables and Data Types}

\subsection{Variable Declaration}

A variable must be declared before it is used.

\begin{lstlisting}[caption={Integer Variable}]
#include <iostream>

int main() {
    int zero = 0;
    std::cout << zero << std::endl;
    return 0;
}
\end{lstlisting}

A declaration tells the compiler:

\begin{itemize}
    \item The variable type
    \item The variable name
    \item The value being stored
\end{itemize}

\subsection{Common Data Types}

\begin{table}[H]
\centering
\caption{Common C++ Data Types}
\begin{tabular}{lll}
\toprule
Type & Example & Description \\
\midrule
\texttt{int} & \texttt{42} & Whole number \\
\texttt{double} & \texttt{3.14159} & Double-precision decimal \\
\texttt{float} & \texttt{2.5f} & Single-precision decimal \\
\texttt{char} & \texttt{'A'} & Single character \\
\texttt{bool} & \texttt{true} & Boolean value \\
\texttt{std::string} & \texttt{"Hello"} & Text string \\
\bottomrule
\end{tabular}
\end{table}

\subsection{String Example}

\begin{lstlisting}[caption={String Variable}]
#include <iostream>
#include <string>

int main() {
    std::string name = "Chabod";
    std::cout << "Name: " << name << std::endl;
    return 0;
}
\end{lstlisting}

% ============================================================
% SECTION VII
% ============================================================
\section{Operators}

\subsection{Arithmetic Operators}

\begin{table}[H]
\centering
\caption{Arithmetic Operators}
\begin{tabular}{lll}
\toprule
Operator & Name & Example \\
\midrule
\texttt{+} & Addition & \texttt{a + b} \\
\texttt{-} & Subtraction & \texttt{a - b} \\
\texttt{*} & Multiplication & \texttt{a * b} \\
\texttt{/} & Division & \texttt{a / b} \\
\texttt{\%} & Modulus & \texttt{a \% b} \\
\bottomrule
\end{tabular}
\end{table}

\subsection{Assignment Operator}

\begin{lstlisting}[caption={Assignment Example}]
int speed = 25;
speed = 30;
\end{lstlisting}

\subsection{Comparison Operators}

\begin{lstlisting}[caption={Comparison Operators}]
int a = 10;
int b = 20;

bool result1 = a == b;
bool result2 = a < b;
bool result3 = a != b;
\end{lstlisting}

% ============================================================
% SECTION VIII
% ============================================================
\section{Control Flow}

\subsection{If Statements}

\begin{lstlisting}[caption={If Statement}]
#include <iostream>

int main() {
    int battery = 75;

    if (battery > 50) {
        std::cout << "Battery level is sufficient." << std::endl;
    } else {
        std::cout << "Battery level is low." << std::endl;
    }

    return 0;
}
\end{lstlisting}

\subsection{For Loop}

\begin{lstlisting}[caption={For Loop}]
#include <iostream>

int main() {
    for (int i = 0; i < 5; i++) {
        std::cout << "Step: " << i << std::endl;
    }

    return 0;
}
\end{lstlisting}

\subsection{While Loop}

\begin{lstlisting}[caption={While Loop}]
#include <iostream>

int main() {
    int altitude = 0;

    while (altitude < 100) {
        altitude += 20;
        std::cout << "Altitude: " << altitude << " m" << std::endl;
    }

    return 0;
}
\end{lstlisting}

% ============================================================
% SECTION IX
% ============================================================
\section{Functions}

\subsection{Basic Function}

\begin{lstlisting}[caption={Basic Function}]
#include <iostream>

void greetUser() {
    std::cout << "Hello from a function." << std::endl;
}

int main() {
    greetUser();
    return 0;
}
\end{lstlisting}

\subsection{Function with Parameters}

\begin{lstlisting}[caption={Function with Parameters}]
#include <iostream>

void printSpeed(int speed) {
    std::cout << "Speed: " << speed << " m/s" << std::endl;
}

int main() {
    printSpeed(45);
    return 0;
}
\end{lstlisting}

\subsection{Function with Return Value}

\begin{lstlisting}[caption={Function with Return Value}]
#include <iostream>

int square(int x) {
    return x * x;
}

int main() {
    int result = square(5);
    std::cout << "Result: " << result << std::endl;
    return 0;
}
\end{lstlisting}

\conceptbox{
A function is a reusable block of code that performs a specific task. It improves structure, readability, and maintainability.
}

% ============================================================
% SECTION X
% ============================================================
\section{Arrays and Vectors}

\subsection{Arrays}

\begin{lstlisting}[caption={Array Example}]
#include <iostream>

int main() {
    int sensorValues[4] = {10, 20, 30, 40};

    std::cout << sensorValues[0] << std::endl;

    return 0;
}
\end{lstlisting}

\subsection{Vectors}

Vectors are dynamic arrays provided by the C++ standard library.

\begin{lstlisting}[caption={Vector Example}]
#include <iostream>
#include <vector>

int main() {
    std::vector<int> values = {10, 20, 30};

    values.push_back(40);

    for (int value : values) {
        std::cout << value << std::endl;
    }

    return 0;
}
\end{lstlisting}

\important{
Arrays have fixed size. Vectors can grow and shrink dynamically, making them more flexible for most modern C++ programs.
}

% ============================================================
% SECTION XI
% ============================================================
\section{Pointers and References}

\subsection{References}

A reference is an alias for another variable.

\begin{lstlisting}[caption={Reference Example}]
#include <iostream>

int main() {
    int speed = 50;
    int& speedRef = speed;

    speedRef = 60;

    std::cout << speed << std::endl;

    return 0;
}
\end{lstlisting}

\subsection{Pointers}

A pointer stores the memory address of another variable.

\begin{lstlisting}[caption={Pointer Example}]
#include <iostream>

int main() {
    int speed = 50;
    int* ptr = &speed;

    std::cout << "Value: " << speed << std::endl;
    std::cout << "Address: " << ptr << std::endl;
    std::cout << "Dereferenced value: " << *ptr << std::endl;

    return 0;
}
\end{lstlisting}

\important{
Pointers are powerful but dangerous if used incorrectly. Invalid pointers, dangling pointers, and memory leaks are common sources of C++ bugs.
}

% ============================================================
% SECTION XII
% ============================================================
\section{Object-Oriented Programming}

\subsection{Classes}

A class groups data and behaviour into one structure.

\begin{lstlisting}[caption={Class Example}]
#include <iostream>
#include <string>

class UAV {
public:
    std::string name;
    double mass;

    void printInfo() {
        std::cout << name << " has mass " << mass << " kg." << std::endl;
    }
};

int main() {
    UAV aircraft;
    aircraft.name = "Test UAV";
    aircraft.mass = 2.5;

    aircraft.printInfo();

    return 0;
}
\end{lstlisting}

\subsection{Constructors}

A constructor initializes an object when it is created.

\begin{lstlisting}[caption={Constructor Example}]
#include <iostream>
#include <string>

class UAV {
public:
    std::string name;
    double mass;

    UAV(std::string n, double m) {
        name = n;
        mass = m;
    }

    void printInfo() {
        std::cout << name << " has mass " << mass << " kg." << std::endl;
    }
};

int main() {
    UAV aircraft("Quadrotor", 1.8);
    aircraft.printInfo();

    return 0;
}
\end{lstlisting}

% ============================================================
% SECTION XIII
% ============================================================
\section{Structs}

\subsection{Struct Example}

Structs are useful for grouping related data.

\begin{lstlisting}[caption={Struct Example}]
#include <iostream>
#include <string>

struct Motor {
    std::string model;
    double thrust;
    double voltage;
};

int main() {
    Motor motor = {"AXI-2212", 12.5, 11.1};

    std::cout << "Motor: " << motor.model << std::endl;
    std::cout << "Thrust: " << motor.thrust << " N" << std::endl;

    return 0;
}
\end{lstlisting}

\conceptbox{
In C++, structs and classes are very similar. The main default difference is that struct members are public by default, while class members are private by default.
}

% ============================================================
% SECTION XIV
% ============================================================
\section{File Input and Output}

\subsection{Writing to a File}

\begin{lstlisting}[caption={Write to File}]
#include <fstream>

int main() {
    std::ofstream file("output.txt");

    file << "C++ file output example." << std::endl;

    file.close();

    return 0;
}
\end{lstlisting}

\subsection{Reading from a File}

\begin{lstlisting}[caption={Read from File}]
#include <iostream>
#include <fstream>
#include <string>

int main() {
    std::ifstream file("output.txt");
    std::string line;

    while (std::getline(file, line)) {
        std::cout << line << std::endl;
    }

    file.close();

    return 0;
}
\end{lstlisting}

% ============================================================
% SECTION XV
% ============================================================
\section{Coding Standards}

\subsection{Why Coding Standards Matter}

Code is read far more often than it is written. Clear formatting, consistent naming, and logical structure make programs easier to understand, debug, and maintain.

Good coding standards include:

\begin{itemize}
    \item Clear variable names
    \item Consistent indentation
    \item Short, focused functions
    \item Meaningful comments
    \item Proper file organisation
    \item Avoiding unnecessary complexity
\end{itemize}

\subsection{Comments}

Single-line comments use:

\begin{lstlisting}[caption={Single-Line Comment}]
// This is a single-line comment
\end{lstlisting}

Multi-line comments use:

\begin{lstlisting}[caption={Multi-Line Comment}]
/*
This is a multi-line comment.
It can span several lines.
*/
\end{lstlisting}

\important{
Comments should explain intent, not repeat obvious code. Good code should be readable first; comments should add context where needed.
}

% ============================================================
% SECTION XVI
% ============================================================
\section{Personal Learning Log}

\subsection{Concepts I Need to Review}

\begin{itemize}
    \item Compilation and linking
    \item Header files
    \item Namespaces
    \item Functions
    \item References
    \item Pointers
    \item Classes
    \item Constructors
    \item Destructors
    \item Memory allocation
    \item Smart pointers
    \item Templates
    \item Standard Template Library
\end{itemize}

\subsection{Mistakes to Avoid}

\begin{itemize}
    \item Forgetting semicolons at the end of statements.
    \item Forgetting to include required header files.
    \item Confusing \texttt{=} with \texttt{==}.
    \item Using variables before declaring them.
    \item Returning from \texttt{main()} before code has executed.
    \item Using pointers without understanding memory addresses.
    \item Forgetting that arrays use zero-based indexing.
    \item Overusing \texttt{using namespace std;} in larger projects.
\end{itemize}

\subsection{Key Mental Model}

\conceptbox{
C++ is not just about syntax. It is about understanding how source code becomes an executable program, how memory is represented, and how software interacts with hardware-level behaviour.
}

\subsection{Questions to Answer Later}

\begin{itemize}
    \item When should I use references instead of pointers?
    \item When should I use a class instead of a struct?
    \item How does dynamic memory allocation work?
    \item What are smart pointers and when should I use them?
    \item How do templates work?
    \item How does C++ support robotics and embedded systems?
    \item How do I structure large C++ projects professionally?
\end{itemize}

% ============================================================
% END DOCUMENT
% ============================================================
\end{document}

